\chapter{Causal Inference and Randomized Controlled Trials}
\label{chap:rct}

\section*{Outline}

This chapter introduces the potential-outcomes language of causal inference, the
average treatment effect, and randomized controlled trials (RCTs).  We then study
two estimators: the difference in means and regression adjustment.  The second
method uses baseline covariates to improve precision without requiring the outcome
to be linear in those covariates.

\section{Potential outcomes}
\label{sec:rct-potential-outcomes}

There is more than one way to formalize causality.  We use the potential-outcomes
framework, associated with Neyman and Rubin.  Imagine an individual indexed by
$i$.

\begin{itemize}
  \item The \emph{treatment} $D_i$ records the treatment that individual $i$
  receives.  Treatments may be discrete or continuous, but here
  $D_i\in\{0,1\}$.
  \item The \emph{potential outcome} $Y_i(d)$ is the outcome that would be
  observed for individual $i$ under treatment status $d$.
  For example, $Y_i(1)$ could be future income if $i$ attends HKUST and
  $Y_i(0)$ future income otherwise.
  \item The observed outcome satisfies the \emph{consistency relation}
  \begin{equation}
    Y_i=Y_i(D_i)=D_iY_i(1)+(1-D_i)Y_i(0).
    \label{eq:rct-consistency}
  \end{equation}
\end{itemize}

Writing $Y_i(d)$ presumes that the treatment is well defined and that one
individual's potential outcome is unaffected by other individuals' treatment
assignments.  Together with consistency, these conditions are commonly grouped
under the \emph{stable unit treatment value assumption} (SUTVA).  SUTVA rules out,
for example, spillovers between treated and untreated students.  Whether it is
credible depends on the application.

The individual treatment effect is
\[
  Y_i(1)-Y_i(0).
\]
It is mathematically well defined, but both potential outcomes cannot normally be
observed for the same individual.  This is the \emph{fundamental problem of causal
inference}.  Replacing the missing $Y_i(0)$ by another person's observed outcome
would require those two people to have the same untreated potential outcome---an
implausibly strong requirement.

\section{The average treatment effect}
\label{sec:rct-ate}

We therefore begin with an average causal parameter:
\begin{equation}
  \tau\equiv \operatorname{ATE}
  =\E\!\left[Y_i(1)-Y_i(0)\right]
  =\E[Y_i(1)]-\E[Y_i(0)].
  \label{eq:rct-ate}
\end{equation}
Under i.i.d. sampling, this population expectation does not depend on the label
$i$.  Notice that treatment effects may still differ across individuals; the ATE
summarizes their population average.

\section{Randomization and identification}
\label{sec:rct-identification}

In an RCT, the researcher assigns treatment at random.  In the simple Bernoulli
design used for the population calculations below,
\begin{equation}
  D_i\indep \bigl(Y_i(1),Y_i(0)\bigr),
  \qquad \Pr(D_i=1)=\pi\in(0,1).
  \label{eq:rct-randomization}
\end{equation}
For example, subjects in a drug trial may be randomly assigned either the drug or
a placebo.  Assignment then does not depend on baseline health, income, or the
outcomes that the subjects would experience under either arm.  A survey that merely
compares people who chose to take the drug with people who did not would not, in
general, have this property.

Randomization turns the unobserved potential-outcome means into observable
conditional means:
\begin{align}
  \tau
  &=\E[Y_i(1)]-\E[Y_i(0)] \notag\\
  &=\E[Y_i(1)\mid D_i=1]-\E[Y_i(0)\mid D_i=0] \notag\\
  &=\E[Y_i\mid D_i=1]-\E[Y_i\mid D_i=0].
  \label{eq:rct-id}
\end{align}
The second line uses random assignment, and the third uses consistency.  Equation
\eqref{eq:rct-id} \emph{identifies} the ATE: it expresses a causal parameter in
terms of the distribution of observed variables.

\section{Difference in means}
\label{sec:rct-dm}

Let $n_d=\sum_{i=1}^n\1\{D_i=d\}$.  The sample analogue of
\eqref{eq:rct-id} is
\begin{equation}
  \widehat\tau_{\mathrm{DM}}
  =\frac{1}{n_1}\sum_{i:D_i=1}Y_i
   -\frac{1}{n_0}\sum_{i:D_i=0}Y_i.
  \label{eq:rct-dm}
\end{equation}
Conditional on having observations in both arms, this estimator is unbiased under
random assignment.  It is also consistent and asymptotically normal.  If treatment
is assigned independently with probability $\pi$, then
\begin{equation}
  \sqrt n\bigl(\widehat\tau_{\mathrm{DM}}-\tau\bigr)
  \dto \N(0,V_{\mathrm{DM}}),
  \qquad
  V_{\mathrm{DM}}
  =\frac{\Var(Y_i(1))}{\pi}
   +\frac{\Var(Y_i(0))}{1-\pi}.
  \label{eq:rct-dm-clt}
\end{equation}
Equivalently, the two variances in \eqref{eq:rct-dm-clt} can be written as
$\Var(Y_i\mid D_i=d)$.  A heteroskedasticity-robust standard error estimates the
sampling uncertainty without imposing equal outcome variances across the two arms.

\subsection{Difference in means as OLS}

Because $D$ is binary,
\begin{align*}
  \E[Y\mid D]
  &=\E[Y\mid D=0]
    +D\{\E[Y\mid D=1]-\E[Y\mid D=0]\}\\
  &=\delta_0+\tau D,
\end{align*}
where the last equality uses \eqref{eq:rct-id}.  Defining
$\varepsilon=Y-\E[Y\mid D]$ gives
\begin{equation}
  Y_i=\delta_0+\tau D_i+\varepsilon_i,
  \qquad \E[\varepsilon_i\mid D_i]=0.
  \label{eq:rct-short-reg}
\end{equation}
This is a population projection identity, not an assumption that the outcome has a
structurally linear model.

The OLS slope from \eqref{eq:rct-short-reg} is exactly
\begin{align*}
  \widehat\tau
  &=\frac{\sum_i(D_i-\bar D)Y_i}
          {\sum_i(D_i-\bar D)^2}\\
  &=\frac{1}{n_1}\sum_{i:D_i=1}Y_i
    -\frac{1}{n_0}\sum_{i:D_i=0}Y_i
   =\widehat\tau_{\mathrm{DM}}.
\end{align*}
Thus a regression of $Y$ on an intercept and $D$ is simply a convenient way to
compute the difference in means and its robust standard error.

\subsection{A survey-wording experiment}
\label{sec:rct-gss}

The following example, adapted from the machine-learning causal-inference tutorial
by Susan Athey and Stefan Wager, uses a simplified General Social Survey experiment.
Roughly half the respondents were asked whether the government spends too much on
``welfare'' ($D_i=1$); the others were asked about ``assistance to the poor''
($D_i=0$).  The binary outcome equals one for a ``yes'' response.  The code verifies
that the direct calculation and OLS produce the same point estimate.

\begin{rcode}
library(lmtest)
library(sandwich)

gss <- read.csv("data/welfare-small.csv")
Y <- gss$Y
D <- gss$D

# Direct difference in means
tau_dm <- mean(Y[D == 1]) - mean(Y[D == 0])
tau_dm

# The coefficient on D is exactly the same number
fit_dm <- lm(Y ~ D)
coeftest(fit_dm, vcov. = vcovHC(fit_dm, type = "HC2"))
\end{rcode}

The wording produces a statistically detectable difference in respondents'
interpretation of government spending.

\section{Regression adjustment}
\label{sec:rct-reg-adjustment}

The difference in means ignores baseline covariates that predict outcomes.  Let
$W_i$ contain such pre-treatment variables and strengthen
\eqref{eq:rct-randomization} to
\begin{equation}
  D_i\indep\bigl(Y_i(1),Y_i(0),W_i\bigr).
  \label{eq:rct-randomization-w}
\end{equation}
Randomization makes covariate adjustment unnecessary for identification, but useful
for precision.

The safest general-purpose specification allows the relationship between $Y$ and
$W$ to differ by treatment arm.  Center $W_i$ at its population mean,
$X_i=W_i-\E[W_i]$, and define the arm-specific linear-projection slopes
\begin{equation}
  \gamma_d
  =\{\E[X_iX_i']\}^{-1}\E\!\left[X_i\{Y_i(d)-\E[Y_i(d)]\}\right],
  \qquad d\in\{0,1\},
  \label{eq:rct-arm-slopes}
\end{equation}
assuming $\E[X_iX_i']$ is nonsingular.  These slopes always exist under the stated
moment conditions.  They do \emph{not} assert that the conditional mean of
$Y_i(d)$ is linear.

In practice, replace $\E[W_i]$ by the overall sample mean $\bar W$ and estimate
the fully interacted regression
\begin{equation}
  Y_i=\alpha+\tau D_i+(W_i-\bar W)'\beta_0
      +D_i(W_i-\bar W)'\bigl(\beta_1-\beta_0\bigr)+u_i.
  \label{eq:rct-lin-reg}
\end{equation}
The coefficient on $D_i$ compares the fitted treatment and control outcomes at the
overall covariate mean.  It consistently estimates the ATE under randomization,
even if the fitted linear regressions are misspecified.  Use a
heteroskedasticity-robust standard error.  The estimator need not be exactly unbiased
in a finite sample, but its bias is asymptotically negligible.

\subsection{Why the interacted adjustment improves precision}

This argument is worth doing carefully.  Define the arm-specific projection
residuals
\[
  r_i(d)=Y_i(d)-\E[Y_i(d)]-X_i'\gamma_d.
\]
The normal equations imply $\E[X_i r_i(d)]=0$, so
\begin{equation}
  \Var(Y_i(d))=\Var(r_i(d))+\gamma_d'\E[X_iX_i']\gamma_d.
  \label{eq:rct-projection-pythagoras}
\end{equation}
Because the regression is evaluated at the overall sample covariate mean, sampling
variation in that mean also enters when the two arm-specific slopes differ.  The
interacted regression-adjustment estimator has asymptotic variance
\begin{equation}
  V_{\mathrm{RA}}
  =\frac{\Var(r_i(1))}{\pi}
   +\frac{\Var(r_i(0))}{1-\pi}
   +(\gamma_1-\gamma_0)'\E[X_iX_i'](\gamma_1-\gamma_0).
  \label{eq:rct-ra-var}
\end{equation}
Combining \eqref{eq:rct-dm-clt}, \eqref{eq:rct-projection-pythagoras}, and
\eqref{eq:rct-ra-var} yields
\begin{equation}
  V_{\mathrm{DM}}-V_{\mathrm{RA}}
  =\left\{
      \sqrt{\frac{1-\pi}{\pi}}\gamma_1
      +\sqrt{\frac{\pi}{1-\pi}}\gamma_0
    \right\}'
    \E[X_iX_i']
    \left\{
      \sqrt{\frac{1-\pi}{\pi}}\gamma_1
      +\sqrt{\frac{\pi}{1-\pi}}\gamma_0
    \right\}
  \geq 0.
  \label{eq:rct-efficiency-gain}
\end{equation}
The gain is large when baseline covariates explain substantial outcome variation.
No invalid equality between treatment-arm conditional covariances is needed: each
arm has its own projection slope and residual variance.

A regression without treatment--covariate interactions is still consistent for the
ATE under \eqref{eq:rct-randomization-w}.  With unequal treatment probabilities and
heterogeneous covariate slopes, however, it does not have the same general
no-loss-of-precision guarantee.  The centered, fully interacted specification in
\eqref{eq:rct-lin-reg} is therefore the preferred default.

For the survey experiment, the implementation is:
\begin{rcode}
covariates <- c("age", "polviews", "income", "educ", "marital", "sex")
gss_c <- gss
gss_c[covariates] <- lapply(gss[covariates], function(x) x - mean(x))

fit_ra <- lm(
  Y ~ D * (age + polviews + income + educ + marital + sex),
  data = gss_c
)
coeftest(fit_ra, vcov. = vcovHC(fit_ra, type = "HC2"))["D", ]
\end{rcode}

Only pre-treatment covariates should be used.  Adjusting for variables caused by
treatment can change the estimand and introduce bias.

\subsection{Monte Carlo illustration}

The next experiment uses a deliberately nonlinear outcome model.  Because $D$ is
randomized independently of $W$, the true ATE is
\[
  \E[1+W^2]=1+\{\Var(W)+\E[W]^2\}=6
\]
when $W\sim\N(2,1)$.  Both estimators are consistent; the interacted adjustment uses
the predictive information in $W$ and is more precise.

\begin{rcode}
set.seed(5280)
n <- 1000
B <- 500
tau_dm <- numeric(B)
tau_ra <- numeric(B)

for (b in seq_len(B)) {
  W <- rnorm(n, 2, 1)
  D <- rbinom(n, 1, 0.5)
  e <- rnorm(n)
  Y <- 1 + D + W^2 * D + exp(W) + e

  tau_dm[b] <- coef(lm(Y ~ D))["D"]
  Wc <- W - mean(W)
  tau_ra[b] <- coef(lm(Y ~ D * Wc))["D"]
}

c(mean_dm = mean(tau_dm), mean_ra = mean(tau_ra))
c(var_dm = var(tau_dm), var_ra = var(tau_ra))
\end{rcode}

\section{Summary}

\begin{itemize}
  \item Potential outcomes define causal effects, while consistency links the
  observed outcome to the potential outcome selected by treatment.
  \item Random assignment identifies the ATE as a difference in observable means.
  \item The difference in means is exactly the coefficient on treatment in a
  regression with an intercept and treatment only.
  \item Baseline covariates are unnecessary for identification in an RCT, but a
  centered, fully interacted regression adjustment can improve precision without
  assuming a correct linear outcome model.
  \item The main assumptions in this chapter are SUTVA, i.i.d. sampling, randomized
  treatment, suitable moments, and nonsingularity of the covariate second-moment
  matrix.  Natural or quasi-experiments may sometimes justify random assignment, but
  that claim must be defended in each application.
\end{itemize}
